% Annexe B -- Probabilites : chaines de Markov finies. A inclure apres \appendix.
% Etiquettes prefixees appB: ; references prefixees app: (refs_app.tex).
\chapter{Chaînes de Markov finies}\label{appB:chap}

Le chapitre~4 étudie trois chaînes : Wright--Fisher, Moran, et la chaîne des configurations sur un graphe. Toutes trois
sont des chaînes de Markov sur un ensemble fini, avec des états absorbants. Cette annexe rassemble, avec leurs
démonstrations, les résultats généraux sur ces chaînes : matrice fondamentale, probabilités et temps moyens
d'absorption, fonctions harmoniques, principe du maximum, martingales et théorème d'arrêt. Le chapitre~4 en démontre des
cas particuliers ; les renvois sont indiqués au fil du texte. Les références classiques sont Kemeny et
Snell~\cite{ch4:KS}, Norris~\cite{app:norris} et Williams~\cite{app:williams}.

\begin{warnbox}[: niveaux]
Les démonstrations de cette annexe sont des démonstrations humaines \tierC. Les cas particuliers formalisés au
chapitre~4 (absorption de Wright--Fisher et de Moran, unicité de la fonction harmonique, théorème isotherme) sont de
niveau \tierA{} dans les fichiers \texttt{formal\_cert/Programmes/PopGen/} et \texttt{EvoGraph/} ; les énoncés généraux
ci-dessous ne le sont pas.
\end{warnbox}

%=====================================================================================================================
\section{Matrices stochastiques et chaînes de Markov}\label{appB:sec:def}

\begin{definition}
Soit $E$ un ensemble fini. Une matrice $P=(P(i,j))_{i,j\in E}$ est \emph{stochastique} si $P(i,j)\ge0$ pour tous $i,j$
et $\sum_jP(i,j)=1$ pour tout $i$. Une suite de variables aléatoires $(X_n)_{n\ge0}$ à valeurs dans $E$ est une
\emph{chaîne de Markov de matrice $P$} si, pour tout $n$ et tous $i_0,\dots,i_{n+1}$ tels que l'événement conditionnant
soit de probabilité non nulle,
\[\mathbb P(X_{n+1}=i_{n+1}\mid X_0=i_0,\dots,X_n=i_n)=P(i_n,i_{n+1}) .\]
\end{definition}
On note $\mathbb P_i$ et $\mathbb E_i$ la probabilité et l'espérance quand $X_0=i$. Une matrice stochastique a la valeur
propre $1$, de vecteur propre $\mathbf 1=(1,\dots,1)^{\mathsf T}$, puisque $P\mathbf 1=\mathbf 1$.

\begin{proposition}[Chapman--Kolmogorov]\label{appB:prop:ck}
Pour tous $n\ge0$ et $i,j\in E$, $\mathbb P_i(X_n=j)=P^n(i,j)$. Plus généralement, si la loi de $X_0$ est le vecteur
ligne $\pi_0$, la loi de $X_n$ est $\pi_0P^n$.
\end{proposition}

\begin{humanproof}
Par récurrence sur $n$. Pour $n=0$, $P^0=I$. Si c'est vrai au rang $n$, la formule des probabilités totales et la
propriété de Markov donnent
$\mathbb P_i(X_{n+1}=j)=\sum_k\mathbb P_i(X_n=k)\,P(k,j)=\sum_kP^n(i,k)P(k,j)=P^{n+1}(i,j)$.
\end{humanproof}

\begin{example}[Les chaînes du chapitre~4]
\begin{itemize}
\item \emph{Wright--Fisher neutre}, $E=\{0,\dots,N\}$ : la loi de $X_{n+1}$ sachant $X_n=i$ est binomiale,
$P(i,j)=\binom Nj(i/N)^j(1-i/N)^{N-j}$.
\item \emph{Moran avec sélection}, $E=\{0,\dots,N\}$ : $P(i,i+1)=p_i$, $P(i,i-1)=q_i$, $P(i,i)=1-p_i-q_i$, avec $p_i=rq_i$
(chapitre~5, §\ref{ch5:U5}). C'est une \emph{chaîne de naissance et de mort}.
\item \emph{Configurations sur un graphe}, $E=\mathcal P(V)$ : l'état est l'ensemble des sommets mutants.
\end{itemize}
Dans les trois cas, les états « tous résidents » et « tous mutants » sont absorbants.
\end{example}

%=====================================================================================================================
\section{États absorbants et matrice fondamentale}\label{appB:sec:abs}

\begin{definition}
Un état $a$ est \emph{absorbant} si $P(a,a)=1$. On note $A$ l'ensemble des états absorbants et $T=E\setminus A$ celui
des autres états, dits \emph{transitoires} dans ce contexte. La chaîne est \emph{absorbante} si $A\ne\emptyset$ et si,
depuis tout état de $T$, un état de $A$ est accessible : pour tout $i\in T$ il existe $n$ et $a\in A$ avec $P^n(i,a)>0$.
\end{definition}

En rangeant les états de $T$ avant ceux de $A$, la matrice prend la \emph{forme canonique}
\[P=\begin{pmatrix}Q&R\\0&I\end{pmatrix},\qquad P^n=\begin{pmatrix}Q^n&(I+Q+\dots+Q^{n-1})R\\0&I\end{pmatrix},\]
où $Q$ est la restriction de $P$ à $T\times T$ et $R$ sa restriction à $T\times A$. La seconde formule se démontre par
récurrence en effectuant le produit par blocs.

\begin{lemma}[Décroissance géométrique]\label{appB:lem:geom}
Si la chaîne est absorbante, il existe $m\ge1$ et $\alpha<1$ tels que $\sum_{j\in T}Q^m(i,j)\le\alpha$ pour tout $i\in
T$. Par conséquent $\sum_{j\in T}Q^n(i,j)\le\alpha^{\lfloor n/m\rfloor}$ et $Q^n\to0$.
\end{lemma}

\begin{humanproof}
Pour chaque $i\in T$, il existe $n_i$ avec $\mathbb P_i(X_{n_i}\in A)>0$. Comme $A$ est absorbant, $\mathbb P_i(X_n\in
A)$ croît avec $n$. Avec $m=\max_in_i$, on a donc $\mathbb P_i(X_m\in A)>0$ pour tout $i\in T$, et
$\sum_{j\in T}Q^m(i,j)=\mathbb P_i(X_m\in T)=1-\mathbb P_i(X_m\in A)$. On pose $\alpha=\max_{i\in T}\bigl(1-\mathbb P_i(X_m\in
A)\bigr)<1$ ($T$ est fini). Notons $s_n=\max_i\sum_jQ^n(i,j)$ ; les coefficients de $Q$ étant positifs,
$s_{n+m}\le s_n\,s_m\le\alpha\,s_n$, et $s_n\le1$.
\end{humanproof}

C'est la version générale du lemme~\ref{ch4:lem:geom} du chapitre~4.

\begin{theorem}[Matrice fondamentale]\label{appB:thm:fond}
Si la chaîne est absorbante, la matrice $I-Q$ est inversible, et
\[N=(I-Q)^{-1}=\sum_{n\ge0}Q^n .\]
Le coefficient $N(i,j)$ est l'espérance du nombre de visites en $j$ partant de $i$ : $N(i,j)=\mathbb
E_i\bigl[\sum_{n\ge0}\mathbf 1_{X_n=j}\bigr]$.
\end{theorem}

\begin{humanproof}
D'après le lemme, la série $\sum_nQ^n$ converge (ses coefficients positifs sont majorés par une série géométrique).
Sa somme $N$ vérifie $(I-Q)N=N(I-Q)=\lim_n(I-Q^{n+1})=I$. Enfin, par Chapman--Kolmogorov et le théorème de
convergence monotone, $\mathbb E_i\bigl[\sum_n\mathbf 1_{X_n=j}\bigr]=\sum_nP^n(i,j)=\sum_nQ^n(i,j)$ pour $i,j\in T$.
\end{humanproof}

\begin{theorem}[Absorption : probabilités et temps moyen]\label{appB:thm:absorption}
Soit une chaîne absorbante et $\tau=\inf\{n\ge0:X_n\in A\}$ le temps d'absorption.
\begin{enumerate}
\item $\tau<\infty$ presque sûrement, et même $\mathbb E_i[\tau]<\infty$.
\item Le vecteur des temps moyens $t(i)=\mathbb E_i[\tau]$, $i\in T$, vaut $t=N\mathbf 1$.
\item La matrice $B(i,a)=\mathbb P_i(X_\tau=a)$, $i\in T$, $a\in A$, vaut $B=NR$.
\item $\lim_{n\to\infty}P^n(i,a)=B(i,a)$ pour $i\in T$, $a\in A$.
\end{enumerate}
\end{theorem}

\begin{humanproof}
(1) $\mathbb P_i(\tau>n)=\mathbb P_i(X_n\in T)=\sum_{j\in T}Q^n(i,j)\le\alpha^{\lfloor n/m\rfloor}$ par le
lemme~\ref{appB:lem:geom}. Donc $\mathbb P_i(\tau=\infty)=0$ et $\mathbb E_i[\tau]=\sum_{n\ge0}\mathbb P_i(\tau>n)<\infty$.

(2) Avant l'absorption, chaque pas est une visite d'un état de $T$ : $\tau=\sum_{n\ge0}\sum_{j\in T}\mathbf 1_{X_n=j}$.
En prenant l'espérance, $t(i)=\sum_{j\in T}N(i,j)$.

(3) L'absorption en $a$ a lieu au pas $n+1$ si $X_n=j\in T$ et $X_{n+1}=a$ :
$\mathbb P_i(X_\tau=a)=\sum_{n\ge0}\sum_{j\in T}Q^n(i,j)R(j,a)=(NR)(i,a)$.

(4) La forme canonique de $P^n$ donne $P^n(i,a)=\bigl((I+\dots+Q^{n-1})R\bigr)(i,a)\to(NR)(i,a)$.
\end{humanproof}

\begin{example}[Moran, $N=3$, $r=2$ : le calcul complet]\label{appB:ex:moran3}
Les probabilités de transition sont $p_1=\frac13$, $q_1=\frac16$, $p_2=\frac4{15}$, $q_2=\frac2{15}$ (chapitre~5,
§\ref{ch5:U5}). Avec $T=\{1,2\}$ et $A=\{0,3\}$,
\[Q=\begin{pmatrix}\frac12&\frac13\\[2pt]\frac2{15}&\frac35\end{pmatrix},\qquad
I-Q=\begin{pmatrix}\frac12&-\frac13\\[2pt]-\frac2{15}&\frac25\end{pmatrix},\qquad\det(I-Q)=\frac15-\frac2{45}=\frac7{45},\]
\[N=\frac{45}7\begin{pmatrix}\frac25&\frac13\\[2pt]\frac2{15}&\frac12\end{pmatrix}=\begin{pmatrix}\frac{18}7&\frac{15}7\\[2pt]
\frac67&\frac{45}{14}\end{pmatrix},\qquad t=N\mathbf 1=\begin{pmatrix}\frac{33}7\\[2pt]\frac{57}{14}\end{pmatrix},\qquad
B(\cdot,3)=N\begin{pmatrix}0\\ \frac4{15}\end{pmatrix}=\begin{pmatrix}\frac47\\[2pt]\frac67\end{pmatrix}.\]
On retrouve les probabilités de fixation $\frac47$ et $\frac67$ de la formule $(1-r^{-i})/(1-r^{-N})$, et l'on obtient
en plus les temps moyens d'absorption : $33/7\approx4{,}71$ pas depuis un mutant, $57/14\approx4{,}07$ depuis deux. Ces
fractions ont été recalculées en arithmétique exacte (\texttt{fractions.Fraction} de Python) pour cette annexe
\tierX{} : $N=\bigl(\begin{smallmatrix}18/7&15/7\\6/7&45/14\end{smallmatrix}\bigr)$, $t=(33/7,\,57/14)$,
$B(\cdot,3)=(4/7,\,6/7)$ et $B(\cdot,0)=(3/7,\,1/7)$.
\end{example}

%=====================================================================================================================
\section{Fonctions harmoniques et principe du maximum}\label{appB:sec:harm}

\begin{definition}
Une fonction $h:E\to\R$ est \emph{harmonique} en $i$ si $h(i)=\sum_jP(i,j)h(j)$, c'est-à-dire si $(Ph)(i)=h(i)$. Elle
est \emph{harmonique sur $T$} si elle l'est en tout point de $T$.
\end{definition}
Autrement dit, $h(i)$ est la moyenne des valeurs de $h$ un pas plus tard. Pour une chaîne de naissance et de mort,
l'harmonicité en $i$ s'écrit $p_i\bigl(h(i+1)-h(i)\bigr)=q_i\bigl(h(i)-h(i-1)\bigr)$ : c'est l'équation du premier pas du
chapitre~5.

\begin{theorem}[Principe du maximum]\label{appB:thm:max}
Soit une chaîne absorbante et $h$ harmonique sur $T$. Alors $\max_Eh=\max_Ah$ et $\min_Eh=\min_Ah$.
\end{theorem}

\begin{humanproof}
Soit $M=\max_Eh$, atteint en un état $i_0$. Si $i_0\in A$, c'est fini. Sinon, $i_0\in T$ et
$M=h(i_0)=\sum_jP(i_0,j)h(j)$ est une moyenne de valeurs $\le M$, à poids positifs de somme $1$ : toutes les valeurs
$h(j)$ avec $P(i_0,j)>0$ valent $M$. De proche en proche, $h=M$ en tout état accessible depuis $i_0$, en particulier
en un état de $A$ (la chaîne est absorbante). Pour le minimum, on applique le résultat à $-h$.
\end{humanproof}

\begin{corollary}[Problème de Dirichlet]\label{appB:cor:dirichlet}
Pour toute fonction $f:A\to\R$, il existe une unique fonction $h$ harmonique sur $T$ et égale à $f$ sur $A$ : c'est
$h(i)=\mathbb E_i[f(X_\tau)]=\sum_{a\in A}B(i,a)f(a)$ pour $i\in T$.
\end{corollary}

\begin{humanproof}
\emph{Unicité.} La différence $d$ de deux solutions est harmonique sur $T$ et nulle sur $A$ ; par le
théorème~\ref{appB:thm:max}, $0\le d\le0$.
\emph{Existence.} Sur $T$, l'équation s'écrit $h_T=Qh_T+Rf$, soit $(I-Q)h_T=Rf$, qui a la solution $h_T=NRf=Bf$
(théorème~\ref{appB:thm:fond}).
\end{humanproof}

En prenant $A=\{0,N\}$, $f(0)=0$ et $f(N)=1$, on obtient : \emph{la probabilité de fixation est l'unique fonction
harmonique à l'intérieur qui vaut $0$ en $0$ et $1$ en $N$}. C'est exactement la structure des preuves du chapitre~4 :
théorème~\ref{ch4:thm:maxprinc} pour Wright--Fisher, où $h(j)=j/N$ ; théorème~\ref{ch4:thm:moran} pour Moran, où
$h(i)=S_i/S_N$ ; proposition~\ref{ch4:prop:harmgraph} pour les graphes. Dans Lean, le pas « la probabilité de
fixation est harmonique » est remplacé par la limite $P^n(i,N)\to h(i)$ (point~4 du théorème~\ref{appB:thm:absorption}),
qui ne demande pas de définir $\tau$.

\begin{remark}[La connexité compte]
Sans l'hypothèse « absorbante », le principe du maximum est faux. Sur $E=\{0,1,2\}$ avec $P(1,1)=1$ et $0,2$
absorbants, toute fonction est harmonique en $1$ ; $h(1)$ est arbitraire et le problème de Dirichlet a une infinité de
solutions. C'est la raison d'être des hypothèses de variance positive (chapitre~4, §\ref{ch4:sec:wf}) et de forte
connexité (lemme~\ref{ch4:lem:conn}) dans les énoncés Lean.
\end{remark}

%=====================================================================================================================
\section{Martingales discrètes}\label{appB:sec:mart}

Nous restons dans le cadre élémentaire d'une chaîne de Markov finie, où toutes les variables sont bornées.

\begin{definition}
Une suite $(M_n)$ de variables aléatoires réelles bornées, avec $M_n$ fonction de $(X_0,\dots,X_n)$, est une
\emph{martingale} (relativement à $(X_n)$) si $\mathbb E[M_{n+1}\mid X_0,\dots,X_n]=M_n$ pour tout $n$. De façon
équivalente : $\mathbb E[(M_{n+1}-M_n)\,Z]=0$ pour toute variable bornée $Z$ fonction de $(X_0,\dots,X_n)$.
\end{definition}

\begin{proposition}\label{appB:prop:harmmart}
Si $h$ est harmonique sur $E$ tout entier (en particulier sur $T$, et constante sur chaque état absorbant), alors
$M_n=h(X_n)$ est une martingale.
\end{proposition}

\begin{humanproof}
Par la propriété de Markov,
\[\mathbb E[h(X_{n+1})\mid X_0,\dots,X_n]=\sum_jP(X_n,j)h(j)=(Ph)(X_n)=h(X_n) .\]
En un état
absorbant $a$, $(Ph)(a)=h(a)$ automatiquement, car $P(a,a)=1$.
\end{humanproof}

\begin{example}
\begin{itemize}
\item \emph{Wright--Fisher.} L'espérance de la loi binomiale $\mathcal B(N,i/N)$ est $i$, donc $h(j)=j$ est harmonique :
le nombre de mutants $X_n$ est une martingale (chapitre~4, §\ref{ch4:sec:wf}).
\item \emph{Moran avec sélection.} $h(i)=r^{-i}$ est harmonique en $0<i<N$. En effet, l'harmonicité s'écrit
$p_i\bigl(h(i+1)-h(i)\bigr)+q_i\bigl(h(i-1)-h(i)\bigr)=0$, et avec $p_i=rq_i$ :
$rq_i(r^{-i-1}-r^{-i})+q_i(r^{-i+1}-r^{-i})=q_ir^{-i}\bigl((1-r)+(r-1)\bigr)=0$. Donc
$r^{-X_n}$ est une martingale. C'est la martingale qu'utilisent Keller et Uğurlu pour les graphes~\cite{ch4:KU}.
\end{itemize}
\end{example}

\begin{theorem}[Théorème d'arrêt]\label{appB:thm:arret}
Soit $(M_n)$ une martingale et $\tau$ un \emph{temps d'arrêt} : pour tout $n$, l'événement $\{\tau=n\}$ ne dépend que de
$(X_0,\dots,X_n)$.
\begin{enumerate}
\item Si $\tau\le K$ presque sûrement pour un entier $K$, alors $\mathbb E[M_\tau]=\mathbb E[M_0]$.
\item Si $\tau<\infty$ presque sûrement et si $|M_n|\le C$ pour tout $n$, alors $\mathbb E[M_\tau]=\mathbb E[M_0]$.
\end{enumerate}
\end{theorem}

\begin{humanproof}
(1) On écrit $M_\tau=M_0+\sum_{k=0}^{K-1}(M_{k+1}-M_k)\mathbf 1_{\tau>k}$. L'événement $\{\tau>k\}$, complémentaire de
$\{\tau\le k\}$, ne dépend que de $(X_0,\dots,X_k)$ ; par définition d'une martingale, chaque terme de la somme est
d'espérance nulle.

(2) Pour $K$ fixé, $\tau\wedge K=\min(\tau,K)$ est un temps d'arrêt borné, donc $\mathbb E[M_{\tau\wedge
K}]=\mathbb E[M_0]$ par (1). Quand $K\to\infty$, $M_{\tau\wedge K}\to M_\tau$ presque sûrement (car $\tau<\infty$) et
$|M_{\tau\wedge K}|\le C$ : la convergence dominée (théorème~\ref{appA:thm:cvd}, sous sa forme probabiliste) conclut.
\end{humanproof}

\begin{corollary}[Ruine du joueur et formule de Moran]\label{appB:cor:ruine}
Pour la chaîne de Moran avec $r\ne1$, partant de $i$, $\mathbb P_i(X_\tau=N)=\dfrac{1-r^{-i}}{1-r^{-N}}$.
\end{corollary}

\begin{humanproof}
La chaîne est absorbante ($p_i>0$ pour $0<i<N$), donc $\tau<\infty$ presque sûrement (théorème~\ref{appB:thm:absorption}).
La martingale $r^{-X_n}$ est bornée par $\max(1,r^{-N})$. Le théorème d'arrêt donne
$r^{-i}=\mathbb E_i[r^{-X_\tau}]=\varphi_i\,r^{-N}+(1-\varphi_i)\cdot1$, avec $\varphi_i=\mathbb P_i(X_\tau=N)$.
On résout : $\varphi_i=(1-r^{-i})/(1-r^{-N})$.
\end{humanproof}
Avec la martingale $X_n$ de Wright--Fisher, le même calcul donne $\varphi_i=i/N$. Ce sont deux démonstrations
probabilistes des formules que le chapitre~4 démontre par l'algèbre linéaire (récurrence sur les différences
$\delta_i$, puis limite de $P^n$).

%=====================================================================================================================
\section{Chaînes irréductibles et loi stationnaire}\label{appB:sec:stat}

Les chaînes absorbantes ne sont pas les seules du livre : la matrice de remplacement $W$ d'un graphe (chapitre~4,
§\ref{ch4:sec:graphes}) est elle-même une matrice stochastique, en général \emph{irréductible}, et le théorème de
Keller--Uğurlu~\cite{ch4:KU} fait intervenir sa loi stationnaire.

\begin{definition}
Une matrice stochastique $P$ est \emph{irréductible} si, pour tous $i,j$, il existe $n\ge0$ avec $P^n(i,j)>0$. Une
\emph{loi stationnaire} est un vecteur ligne $\pi\ge0$, de somme $1$, tel que $\pi P=\pi$.
\end{definition}

\begin{proposition}[Existence]\label{appB:prop:existence}
Toute matrice stochastique sur un ensemble fini a au moins une loi stationnaire.
\end{proposition}

\begin{humanproof}
Partons d'une loi $\pi_0$ quelconque et posons $\mu_n=\frac1n\sum_{k=0}^{n-1}\pi_0P^k$. Ce sont des lois de
probabilité ; elles vivent dans le simplexe, qui est compact. Une sous-suite $(\mu_{n_l})$ converge vers une loi $\mu$.
Or $\mu_nP-\mu_n=\frac1n(\pi_0P^n-\pi_0)$, dont chaque coefficient est majoré par $1/n$ en valeur absolue. À la limite,
$\mu P=\mu$.
\end{humanproof}

\begin{proposition}[Unicité et positivité]\label{appB:prop:unicite}
Si $P$ est irréductible, la loi stationnaire $\pi$ est unique et $\pi(j)>0$ pour tout $j$.
\end{proposition}

\begin{humanproof}
\emph{Les fonctions harmoniques sont constantes.} Si $Ph=h$ sur tout $E$, l'argument du théorème~\ref{appB:thm:max}
montre que $h$ vaut son maximum en tout état accessible depuis un point où il est atteint, donc partout, par
irréductibilité. Le noyau de $P-I$ est donc la droite engendrée par $\mathbf 1$.
\emph{Unicité.} Le noyau à gauche, $\{x:x(P-I)=0\}$, a la même dimension que le noyau à droite (rang d'une matrice et
de sa transposée), donc il est de dimension $1$ ; une seule de ses droites a une somme égale à $1$.
\emph{Positivité.} Si $\pi(i)>0$ (il en existe un) et $P^n(i,j)>0$, alors $\pi(j)=(\pi P^n)(j)\ge\pi(i)P^n(i,j)>0$.
\end{humanproof}

\begin{theorem}[Convergence géométrique]\label{appB:thm:doeblin}
Supposons qu'il existe $m\ge1$ et $\delta>0$ tels que $P^m(i,j)\ge\delta$ pour tous $i,j$ (matrice \emph{primitive}). Soit
$n=|E|$. Pour toutes lois $\mu,\nu$, $\|\mu P^m-\nu P^m\|_1\le(1-n\delta)\|\mu-\nu\|_1$. En particulier $\pi_0P^k\to\pi$
pour toute loi initiale $\pi_0$, à vitesse géométrique.
\end{theorem}

\begin{humanproof}
On a $n\delta\le1$, car une ligne de $P^m$ a $n$ coefficients $\ge\delta$ de somme $1$. Si $n\delta=1$, toutes les lignes
de $P^m$ sont uniformes et $\mu P^m=\nu P^m$. Sinon, écrivons $P^m=n\delta\,U+(1-n\delta)R$, où $U$ a toutes ses
lignes égales à $(\frac1n,\dots,\frac1n)$ ; la matrice $R=(P^m-n\delta U)/(1-n\delta)$ est alors à coefficients
positifs, de lignes de somme $1$. Le vecteur $x=\mu-\nu$ est de somme nulle, donc $xU=0$, et
$xP^m=(1-n\delta)\,xR$. Enfin $\|xR\|_1\le\sum_i|x_i|\sum_jR(i,j)=\|x\|_1$. On applique l'inégalité à $\mu=\pi_0P^k$ et
$\nu=\pi$ (qui vérifie $\pi P^m=\pi$) ; les puissances intermédiaires ne font pas croître la norme, par le même calcul
avec $P$ à la place de $R$.
\end{humanproof}

\begin{example}[Le graphe en étoile à trois sommets]\label{appB:ex:etoile}
Prenons le centre $0$ et les feuilles $1$, $2$, avec $W(0,1)=W(0,2)=\frac12$ et $W(1,0)=W(2,0)=1$ (exercice~\ref{ch4:ex:star}).
L'équation $\pi W=\pi$ donne $\pi(0)=\pi(1)+\pi(2)$ et $\pi(1)=\pi(2)=\pi(0)/2$, d'où $\pi=(\frac12,\frac14,\frac14)$.
Cette matrice est irréductible mais \emph{pas} primitive : elle est périodique de période $2$ (on alterne entre le
centre et les feuilles), et $\pi_0W^k$ ne converge pas si l'on part du centre. Le théorème~\ref{appB:thm:doeblin} ne
s'applique donc pas ; la loi stationnaire existe et est unique quand même. L'étoile n'est pas isotherme ($W$ n'est pas
doublement stochastique : la colonne du centre somme à $2$), ce qui explique que sa probabilité de fixation diffère de
celle de Moran (chapitre~4). Avec des poids de sélection stationnaires au sens de l'énoncé Lean \lean{stationary_fixation}
(\texttt{EvoGraph/BirthDeath.lean}), l'enregistrement N4~\cite{ch4:N4} retrouve la valeur de Moran $4/7$ depuis chacun
des trois sommets (calcul en fractions exactes, \tierX).
\end{example}

\begin{remark}[Et les chaînes absorbantes ?]
Une chaîne absorbante n'est pas irréductible dès qu'elle a un état transitoire. Ses lois stationnaires sont les
mélanges de masses de Dirac sur les états absorbants, et la loi limite dépend du point de départ : c'est la matrice
$B$ du théorème~\ref{appB:thm:absorption}.
\end{remark}

%=====================================================================================================================
\section{Du temps discret au temps continu : l'uniformisation}\label{appB:sec:unif}

La crate \texttt{qf-popgen} (chapitre~5, cas U5) n'itère pas $P^n$ : elle intègre l'équation différentielle linéaire
$\dot\pi(t)=\pi(t)(P-I)$ avec CVODE. Voici pourquoi les deux calculs ont la même limite.

\begin{proposition}[Mélange de Poisson]\label{appB:prop:unif}
La solution de $\dot\pi=\pi(P-I)$, $\pi(0)=\pi_0$, est
\[\pi(t)=\pi_0e^{t(P-I)}=\sum_{n\ge0}e^{-t}\frac{t^n}{n!}\,\pi_0P^n .\]
Si $P$ est absorbante, $\pi(t)(a)\to\sum_i\pi_0(i)B(i,a)$ quand $t\to\infty$, pour tout $a\in A$.
\end{proposition}

\begin{humanproof}
Les matrices $tP$ et $-tI$ commutent, donc
\[e^{t(P-I)}=e^{-t}e^{tP}=e^{-t}\sum_{n\ge0}\frac{t^nP^n}{n!} .\]
Pour la limite, notons
$u_n=\pi_0P^n(a)$, qui converge vers $\ell=\sum_i\pi_0(i)B(i,a)$ (théorème~\ref{appB:thm:absorption}, point~4, et
$P^n(a,a)=1$). Les poids $e^{-t}t^n/n!$ forment une loi de Poisson de paramètre $t$, dont la masse sur $\{n\le n_0\}$
tend vers $0$ quand $t\to\infty$, pour tout $n_0$ fixé. On coupe la somme en $n\le n_0$ et $n>n_0$, avec $n_0$ tel que
$|u_n-\ell|\le\varepsilon$ pour $n>n_0$ ; comme $|u_n-\ell|\le2$ pour tout $n$,
\[|\pi(t)(a)-\ell|\le2\,\mathbb P\bigl(\mathrm{Poisson}(t)\le n_0\bigr)+\varepsilon .\]
\end{humanproof}
Cette proposition est le pas « l'équation différentielle et le noyau à $n$ pas ont la même limite » que le chapitre~5
signale comme non formalisé : elle est de niveau \tierC{} dans ce livre.

%=====================================================================================================================
\section{Exercices}\label{appB:sec:exos}

\begin{exercise}
Pour la chaîne de Moran neutre ($r=1$) à $N$ individus, montrer que $\delta_i=\varphi_{i+1}-\varphi_i$ est constant et
retrouver $\varphi_i=i/N$. Calculer la matrice fondamentale pour $N=3$ et le temps moyen d'absorption depuis $i=1$.
\end{exercise}

\begin{exercise}
Montrer que, pour une chaîne de naissance et de mort absorbante en $0$ et $N$, le vecteur $t$ des temps moyens vérifie
$t(i)=1+p_it(i+1)+q_it(i-1)+(1-p_i-q_i)t(i)$ pour $0<i<N$, avec $t(0)=t(N)=0$. Vérifier sur l'exemple~\ref{appB:ex:moran3}
que $(33/7,57/14)$ est solution.
\end{exercise}

\begin{exercise}
Soit $(X_n)$ la chaîne de Wright--Fisher neutre. Montrer que $X_n(N-X_n)\,(1-1/N)^{-n}$ est une martingale. En déduire
que $\mathbb E_i[X_n(N-X_n)]=i(N-i)(1-1/N)^n$, puis une majoration de $\mathbb P_i(\tau>n)$.
\end{exercise}

\begin{exercise}[Le témoin négatif du théorème d'arrêt]
Sur $\Z$, la marche simple symétrique $S_n$ partant de $0$ est une martingale, et $\tau=\inf\{n:S_n=1\}$ est fini
presque sûrement (on l'admet). Montrer que $\mathbb E[S_\tau]=1\ne\mathbb E[S_0]$. Quelle hypothèse du
théorème~\ref{appB:thm:arret} est violée ?
\end{exercise}
